\chapter{Start here}\label{ch:start}

You are in Claude Desktop. This chapter goes once around the whole loop: get LeanFM, connect
Claude to it, and design a small system with it. The rest of the book explains what happened.

\section{Get LeanFM and start it}

In a terminal, on Linux or macOS:
\begin{lstlisting}[language={},basicstyle=\ttfamily\footnotesize]
git clone https://github.com/rfielding/leanFM.git
cd leanFM
./setup.sh
make build
make serve
\end{lstlisting}
\texttt{./setup.sh} installs only what is missing, and asks for your password once if it needs
system packages. Run it again whenever you like: when nothing is missing it says so. The first
\texttt{make build} takes a few minutes, while Lean builds LeanFM. \texttt{make serve} keeps
running: it is the LeanFM server, answering on \texttt{http://127.0.0.1:8080}. Leave it open.

To build this book as well, run \texttt{./setup.sh book} first; \texttt{./setup.sh all} sets up
everything. With no \texttt{git} yet, install it first: \texttt{sudo apt install git} on Ubuntu,
\texttt{xcode-select --install} on macOS.

\section{Connect Claude}

In Claude Desktop's Code tab, open the \texttt{leanFM} folder. The folder's
\texttt{.mcp.json} connects Claude to the server you just started. For a Claude Desktop chat
instead, run this once in the folder and restart Claude Desktop once:
\begin{lstlisting}[language={},basicstyle=\ttfamily\footnotesize]
python3 scripts/claude_desktop_config.py --root "$PWD" \
  --forward http://127.0.0.1:8080/mcp --install
\end{lstlisting}
From then on, after a rebuild you restart only \texttt{make serve}. Claude Desktop and the
server may start in either order: a request that arrives first waits for the server to listen.

Now ask Claude:
\begin{quote}
\emph{Find the LeanFM MCP server from this directory.}
\end{quote}
Claude lists LeanFM's tools; \texttt{server\_info} names the commit the server was built from.

\section{Design TicTacToe}

Ask:
\begin{quote}
\emph{Let's design a TicTacToe game with LeanFM.}
\end{quote}
Claude writes the game as a protocol: who sends what to whom, and in what order. Two players
and a referee. X starts a game, the referee tells O it has joined, and then the players take
turns through the referee until it announces the result:
\begin{lstlisting}[language={},basicstyle=\ttfamily\footnotesize,breaklines=true]
protocol TicTacToe (X, O, Referee) {
  g: NewGame from X to Referee;
  Joined from Referee to O cause g;
  loop Round {
    x: MoveX from X to Referee;
    choice {
      OTurn from Referee to O cause x;
      o: MoveO from O to Referee;
      choice {
        XTurn from Referee to X cause o;
        continue Round;
      } or {
        Result from Referee to X cause o;
        Result from Referee to O cause o;
      }
    } or {
      Result from Referee to X cause x;
      Result from Referee to O cause x;
    }
  }
}
\end{lstlisting}

LeanFM checks the text and draws it (Figure~\ref{fig:ttt}).

\begin{figure}[ht]
\centering
\AccessibleFigure{Interaction diagram of TicTacToe. NewGame from X to Referee then Joined from Referee to O. Then a loop named Round. MoveX from X to Referee. Then either OTurn from Referee to O and MoveO from O to Referee and then either XTurn from Referee to X and the loop continues or Result from Referee to X and to O. Or Result from Referee to X and to O.}{\ThemedDiagram{width=.62\textwidth}{ctl-tictactoe}}
\caption{The first draft of TicTacToe, as LeanFM draws it.}
\label{fig:ttt}
\end{figure}

Then Claude asks the question every game must answer: does it always end? LeanFM says no, and
shows the run that does not:
\ToolOutput{generated/ctl-tictactoe-draft-ends.tex}
The players can trade moves forever. Nothing in the draft says the board fills up. But it
does: nine squares, X first, so X moves at most five times. One phrase says that, and one
block says what happens when a fifth round would go on:
\begin{lstlisting}[language={},basicstyle=\ttfamily\footnotesize,breaklines=true]
protocol TicTacToe (X, O, Referee) {
  g: NewGame from X to Referee;
  Joined from Referee to O cause g;
  loop Round at most 5 {
    x: MoveX from X to Referee;
    choice {
      OTurn from Referee to O cause x;
      o: MoveO from O to Referee;
      choice {
        XTurn from Referee to X cause o;
        continue Round;
      } or {
        Result from Referee to X cause o;
        Result from Referee to O cause o;
      }
    } or {
      Result from Referee to X cause x;
      Result from Referee to O cause x;
    }
  } otherwise {
    Result from Referee to X;
    Result from Referee to O;
  }
}
\end{lstlisting}

Now every requirement Claude wrote down is decided:
\begin{center}
\ToolOutput{generated/ctl-tictactoe-checks.tex}
\end{center}

\section{What just happened}

That was Requirements Logic at work. The protocol is the design, written so a machine can
read it. Each requirement is a claim about what may and must happen next (a formula of CTL,
the half of the logic Chapter~\ref{ch:checker} explains). LeanFM decided each
one over every run the protocol allows. When one failed, it did not just say so: it showed the
run, and that run was the design mistake.

None of it had to be trusted. Claude wrote the protocol and the formulas; LeanFM checked them.
Part~\ref{part:language} is the protocol language; Part~\ref{part:ctl} is how a protocol
becomes claims and how claims are checked; Part~\ref{part:lists} is the list of requirements a project
keeps; and Part~\ref{part:mcp} is the server Claude just used.
